% Merged: Architecture Patterns (03) + Modalities/Architectures (04)
% Created: 2026-01-27
\subsection{Architectures and Personalization}
\label{sec:architecture-patterns}

\subsubsection{Learning Paradigms}

Deep learning approaches in wearable seizure monitoring employ diverse learning paradigms. Feature-based approaches appear in four studies. These methods extract handcrafted features as input to classifiers. \textcite{fineDetectionKeyAutomated2025} extracted 594 features from 6-axis ACC and gyroscope data for an ANN classifier. The large feature set raises overfitting concerns given the small sample size (n=15), and no study directly compares feature-based to end-to-end approaches. \textcite{wangEpilepticSeizureDetection2025} engineered attitude angle features for an LSTM network with 40 hidden units. \textcite{singhrathoreDevelopmentMultimodalMachine2024} used multi-modal features for an MLP classifier in a single-patient case study.

End-to-end learning appears in two studies, both in forecasting. \textcite{meiselMachineLearningWristband2020} used an LSTM with 10 units operating directly on raw multi-modal wrist data. \textcite{nasseriAmbulatorySeizureForecasting2021} used a 4-layer LSTM with 128 hidden units per layer receiving FFT channels of raw sensor data.

Hybrid approaches appear in two studies. \textcite{vielufSeizureMonitoringCombined2025} combined a DNN with handcrafted harmonic features extracted from 24-hour patterns. \textcite{elemamAutomatedValidatedTool2025} used separate CNNs for HRV and audio data without fusion between modalities.

Ensemble methods appear in detection applications. \textcite{spahrDeepLearningbasedDetection2025} trained 30 separate 1D CNN models with quantile aggregation, achieving 96\% sensitivity with 0.0054/h FPR. The ensemble aggregation quantile serves as a tunable parameter that allows adapting the sensitivity-FPR trade-off according to patient needs. The model (Episave) operates on 3D-ACC data at 32~Hz using 30-second windows. \textcite{dongTwoLayerEnsembleMethod2022} used a two-stage architecture with threshold-based pre-screening followed by a CNN-LSTM model.

Anomaly detection appears in two studies. \textcite{reintjesECGBasedDetectionEpileptic2025} compared three unsupervised anomaly detection methods (Matrix Profile, MADRID, TimeVQVAE) on single-lead ECG data. \textcite{odeDevelopmentEpilepticSeizure2023} used a self-attentive autoencoder for patient-specific anomaly detection on ECG RRI intervals.

\subsubsection{Architecture Patterns by Application}

Deep learning architectures differ substantially between detection and forecasting applications. Forecasting studies concentrate on LSTM-based approaches while detection research is more diverse but CNN-based approaches dominate.

CNN-based methods are prominent in detection. \textcite{spahrDeepLearningbasedDetection2025} implemented an ensemble of 30 CNN models with 14 convolutional layers each, operating on ACC data sampled at 32~Hz with 30-second windows. \textcite{elemamAutomatedValidatedTool2025} used separate CNNs for HRV and audio processing. Traditional neural networks appear in earlier studies, with \textcite{fineDetectionKeyAutomated2025} using an ANN with 594 handcrafted features achieving 100\% sensitivity.

Three of five forecasting studies use LSTM variants. \textcite{meiselMachineLearningWristband2020} implemented a minimal LSTM with 10 hidden units operating directly on raw multi-modal wrist data. \textcite{nasseriAmbulatorySeizureForecasting2021} used a deeper 4-layer LSTM with 128 hidden units per layer. \textcite{stirlingForecastingSeizureLikelihood2021} combined LSTM with random forest and logistic regression in an ensemble. No forecasting study uses CNN-based approaches, possibly because forecasting requires modeling temporal dependencies rather than spatial patterns.

Handcrafted features remain competitive. Three of eight detection studies use feature-based approaches with strong performance. \textcite{fineDetectionKeyAutomated2025} achieved 100\% sensitivity using 594 handcrafted features. This suggests that domain knowledge remains valuable despite the end-to-end learning trend.

Window sizes vary substantially between applications. Detection studies use shorter windows ranging from 4 seconds to 5 minutes. Forecasting studies use longer windows ranging from 30 seconds to 60 minutes.

\subsubsection{Personalization Strategies}

Personalization strategy represents a fundamental design choice. Global models train on population data and apply to all patients while patient-specific models train on individual patient data. Five studies use global model approaches, four studies use patient-specific approaches, and one study uses a mixed approach.

\paragraph{Global Models.}
Five studies use global model approaches. \textcite{spahrDeepLearningbasedDetection2025} trained a global model on 384 patients with tunable parameters. \textcite{fineDetectionKeyAutomated2025} trained on 15 patients and tested on 3 patients. \textcite{dongTwoLayerEnsembleMethod2022} trained on 68 patients with 10-fold cross-validation. \textcite{elemamAutomatedValidatedTool2025} trained on 198 questionnaire responses and 30 HRV recordings. \textcite{meiselMachineLearningWristband2020} trained a global model tested with LOSO validation.

Global models enable on-device processing without requiring individual patient training. \textcite{spahrDeepLearningbasedDetection2025} achieved 112 ms inference time suitable for real-time processing. However, global models may not capture individual seizure patterns. \textcite{meiselMachineLearningWristband2020} used LOSO validation to assess global model generalizability, achieving 62\% patient success (43 of 69 patients above chance).

\paragraph{Patient-Specific Models.}
Four studies use patient-specific approaches. \textcite{stirlingForecastingSeizureLikelihood2021} trained individual models with weekly retraining for 11 patients. \textcite{nasseriAmbulatorySeizureForecasting2021} trained separate models for each of 6 patients. \textcite{odeDevelopmentEpilepticSeizure2023} trained anomaly detection models at the 99\% confidence level for each of 66 patients. \textcite{singhrathoreDevelopmentMultimodalMachine2024} presented a single-patient case study.

Patient-specific models achieve higher individual success rates. \textcite{stirlingForecastingSeizureLikelihood2021} achieved 100\% patient success for hourly forecasting. \textcite{nasseriAmbulatorySeizureForecasting2021} achieved 83\% patient success (5 of 6 patients). These rates are substantially higher than the 62\% achieved by global models with LOSO validation. However, patient-specific approaches require a mean of 14.6 months of data per patient or 60+ days per patient. This data requirement creates a cold start problem where new patients must wait weeks or months before the system becomes useful.

\paragraph{Mixed Approaches.}
\textcite{vielufSeizureMonitoringCombined2025} used a mixed approach. The study combined population-level 24-hour harmonic patterns with individual diary data. This hybrid achieved 82\% sensitivity and 67\% specificity. The approach may balance the advantages of global and patient-specific methods.

\subsubsection{Validation Approaches}

Validation approaches vary across studies. \textcite{reintjesECGBasedDetectionEpileptic2025} used subject-split validation, training on one set of patients and testing on different patients to assess generalization. The results showed poor generalization with FPR ranging from 0.11/h to 65.62/h depending on method and configuration.

\textcite{nasseriAmbulatorySeizureForecasting2021} used temporal split validation. Early data from each patient served as training while later data served as testing. This approach assesses temporal stability rather than generalization to new patients.

\textcite{meiselMachineLearningWristband2020} used LOSO validation to assess global model performance on unseen patients, achieving 62\% patient success.

\subsubsection{Temporal Evolution}

Architecture choices have evolved over time. The earliest study, \textcite{borujenyDetectionEpilepticSeizure2013}, used traditional ANN approaches on MICAz accelerometer data in a 3-patient study. Studies published between 2020 and 2022 focused on LSTM-based forecasting architectures. Studies published between 2023 and 2026 employed more complex architectures including ensembles, hybrid methods, and attention mechanisms.

\subsubsection{Computational Considerations}

Global models enable on-device deployment. \textcite{spahrDeepLearningbasedDetection2025} achieved 112 ms inference time suitable for real-time processing. On-device processing enables continuous monitoring without cloud connectivity.

Patient-specific models often require cloud or smartphone-based processing. \textcite{stirlingForecastingSeizureLikelihood2021} used smartphone-based processing with weekly retraining. \textcite{nasseriAmbulatorySeizureForecasting2021} used cloud-based processing. Cloud-based approaches may experience connectivity issues and latency.

\textbf{Key findings.} Detection research employs architectural diversity including CNNs, LSTMs, ANNs, and anomaly detection. Forecasting research concentrates heavily on LSTM variants with no CNN-based approaches, possibly because forecasting requires modeling temporal dependencies rather than spatial patterns. Patient-specific models achieve higher individual success (83-100\%) compared to global models with LOSO validation (62\%), but they require weeks to months of individual data. Global models enable immediate deployment and on-device processing but may not work well for all patients.

\textbf{Limitations.} Architecture descriptions vary in detail across studies. Some studies report complete architectural details while others provide minimal information. This variability complicates direct comparison and replication. Detection studies rarely report patient-level success. Only two of eight detection studies provide individual patient results. This reporting gap prevents assessment of which patients benefit from global detection models.
